%----------------------------------------------------------------------------------------
%	ABSTRACT PAGES
%----------------------------------------------------------------------------------------

% IMPORTANT NOTE: the abstract must always be written in two languages. If the dissertation
% is written in Portuguese you have selected 'portuguese' as the language in the document class.
% Therefore, the portuguese version of the abstract must come first, so write it in the
% below area denoted by 'MAIN LANGUAGE ABSTRACT'. The english version follows in the
% 'SECOND LANGUAGE ABSTRACT' section.
% If the dissertation is written in English, first will come the abstract in English
% ('MAIN LANGUAGE ABSTRACT') and then in Portuguese ('SECOND LANGUAGE ABSTRACT').

\begin{abstract}
%%%%%%%%%%%%%%%%%%%%%%%%%%%%%% MAIN LANGUAGE ABSTRACT %%%%%%%%%%%%%%%%%%%%%%%%%%%%%%%%%%
A \textit{framework JavaFX}, será estudada e implementada neste relatório técnico-científico, com a criação de uma aplicação distribuída esta tecnologia será apresentada.



Aqui deverá ser apresentado o resumo de todo o trabalho efetuado. Esta secção não deverá exceder uma página.

Deve contextualizar o problema que pretende resolver ou a hipótese que irá formular, procure evidenciar as vantagens e desvantagens (se as houver) da solução encontrada, como também a forma através da qual a solução/hipótese foi validada. Neste último ponto, deverá referir-se aos desenvolvimentos efetuados, e à forma como validou (conformidade) e avaliou (desempenho) a solução encontrada.

A dissertação deve conter sempre duas versões do resumo: uma primeira no idioma do texto principal e a segunda num outro idioma. Este \textit{template} assume que os dois idiomas em consideração são sempre Português e Inglês, assim, a classe irá colocar os cabeçalhos respetivos de acordo com o idioma selecionado nas opções da classe no ficheiro \file{main.tex}. 

%----------------------------------------------------------------------------------------

\vspace*{10mm} 
\noindent
\textbf{\keywordslabel}: \textit{JavaFX}, Lista, separada por vírgulas, de palavras, frases, ou acrónimos chave no âmbito do trabalho descrito neste texto. 

%%%%%%%%%%%%%%%%%%%%%%%%% END OF THE MAIN LANGUAGE ABSTRACT %%%%%%%%%%%%%%%%%%%%%%%%%%%%%%
\end{abstract}
\begin{secondlangabstract}
%%%%%%%%%%%%%%%%%%%%%%%%%%%%%% SECOND LANGUAGE ABSTRACT %%%%%%%%%%%%%%%%%%%%%%%%%%%%%%%%%%
The framework JavaFX, will be studied and implemented in this theoretical-cientific paper, with the creation of distributed aplication, this tecnology will be presented.


%----------------------------------------------------------------------------------------

\vspace*{10mm} 
\noindent
\textbf{\keywordslabel}: JavaFX,

%%%%%%%%%%%%%%%%%%%%%%%%%% END OF THE SECOND LANGUAGE ABSTRACT %%%%%%%%%%%%%%%%%%%%%%%%%%%%%
\end{secondlangabstract}

